\section{Terminal injectivity and the M\"obius boundary}
\label{sec:tpc39-terminal}

The terminal factorization \(s=1\) displays the arithmetic content that
the twist bank must estimate.  We first remove a possible combinatorial
ambiguity: distinct row--alias atoms do not represent the same integer.

\subsection{Exact terminal injectivity}

For fixed \(j\in\cJ\), define
\begin{equation}
 n_{\alpha,j,k}
 :=m_\alpha j+h+kM,
 \qquad -L\le k\le L.
 \label{eq:tpc39-terminal-target}
\end{equation}

\begin{theorem}[Terminal row--alias injectivity]
\label{thm:tpc39-terminal-injectivity}
Assume \eqref{eq:tpc39-interface-prime-range} and, as at the endpoint,
\begin{equation}
 \max_{\alpha,\beta}|m_\alpha-m_\beta|<M.
 \label{eq:tpc39-row-difference-less-M}
\end{equation}
For every fixed \(j\), the map
\begin{equation}
 (\alpha,k)\longmapsto n_{\alpha,j,k}
 \label{eq:tpc39-terminal-injective-map}
\end{equation}
is injective.
\end{theorem}

\begin{proof}
If \(n_{\alpha,j,k}=n_{\beta,j,k'}\), then
\begin{equation}
 (m_\alpha-m_\beta)j=(k'-k)M.
 \label{eq:tpc39-terminal-collision-equation}
\end{equation}
Every prime factor \(q_i\) of \(M\) is larger than the positive integer
\(j\), so \((j,M)=1\).  Thus
\(M\mid m_\alpha-m_\beta\).  By
\eqref{eq:tpc39-row-difference-less-M}, the two row integers are equal.
Source separation gives \(\alpha=\beta\), and then
\eqref{eq:tpc39-terminal-collision-equation} gives \(k=k'\).
\end{proof}

At the endpoint, \(m_\alpha\asymp Q\) and \(M/Q\to\infty\), so
\eqref{eq:tpc39-row-difference-less-M} is automatic after the dyadic
constants are fixed.  The theorem says that no hidden target multiplicity
is responsible for the terminal difficulty.

\subsection{The atomic diagonal is already at the target scale}

With a terminal cutoff \(\Phi_0\), put
\begin{equation}
 b_{\alpha,j,k}
 :=-\mu(n_{\alpha,j,k})\log(n_{\alpha,j,k})
 \Phi_0(n_{\alpha,j,k}),
 \label{eq:tpc39-terminal-b}
\end{equation}
interpreting the term as zero outside the positive terminal support.  A
left terminal slice, including the equality alias \(k=0\), has the form
\begin{equation}
 \cT_{\gamma,j}^{L}
 =\sum_\alpha\sum_{k\in\cI_B}
 c_{\alpha,\gamma,j}^{L}b_{\alpha,j,k},
 \label{eq:tpc39-terminal-left-slice}
\end{equation}
where
\begin{equation}
 c_{\alpha,\gamma,j}^{L}
 =\gamma_\alpha^{(1)}
 \mathfrak A_{\alpha,\gamma}^{\rm act}(j).
 \label{eq:tpc39-terminal-left-coefficient}
\end{equation}
The right slice is symmetric.  The shift changes the terminal atom \(b\),
not the inherited row--time multiplier, and the deterministic multiplier
ledger gives
\begin{equation}
 \sup_{\alpha,\gamma,j}
 \left|\mathfrak A_{\alpha,\gamma}^{\rm act}(j)\right|
 \ll X^{o(1)}.
 \label{eq:tpc39-terminal-multiplier-ledger}
\end{equation}
This is only a pointwise mask, residue, and smooth--weight bound; it is not
the unproved moving--shift Hilbert output estimate.  All cutoffs and
logarithms are absorbed into \(X^{o(1)}\).

For the \(\nu\)-th twist--bank member, each summand additionally carries
\(\omega^{\nu k}\).  This factor has modulus one and disappears on the
atomic diagonal, so the same diagonal bound holds after averaging over
\(\nu\).

Expanding the two squared energies gives
\begin{equation}
 \|\cT^{L}\|_2^2+\|\cT^{R}\|_2^2
 =\cD_{\rm term}+\cO_{\rm term},
 \label{eq:tpc39-terminal-diagonal-offdiagonal}
\end{equation}
where \(\cD_{\rm term}\) contains only
\((\alpha,k)=(\beta,k')\), and \(\cO_{\rm term}\) contains every other
pair.

\begin{proposition}[Terminal diagonal closure]
\label{prop:tpc39-terminal-diagonal-closure}
The atomic diagonal satisfies
\begin{equation}
 \boxed{
 \cD_{\rm term}
 \ll X^{o(1)}Q^2JK_B.}
 \label{eq:tpc39-terminal-diagonal-bound}
\end{equation}
For \(B\asymp X\),
\begin{equation}
 \cD_{\rm term}
 \ll X^{o(1)}\frac{Q^3}{J}.
 \label{eq:tpc39-terminal-diagonal-endpoint}
\end{equation}
\end{proposition}

\begin{proof}
There are \(O(X^{o(1)}Q)\) opposite rows and \(O(J)\) orbit times.
For each of the \(K_B\) aliases, use
\(\sum_\alpha|\gamma_\alpha^{(i)}|^2\ll X^{o(1)}Q\) and the bounded
multiplier.  This proves
\eqref{eq:tpc39-terminal-diagonal-bound} for both sides.  Finally,
\eqref{eq:tpc39-interface-key-ledger} gives
\eqref{eq:tpc39-terminal-diagonal-endpoint}.
\end{proof}

Together with \cref{thm:tpc39-terminal-injectivity}, this closes the
terminal collision and atomic diagonal ledger.  The remaining term
\(\cO_{\rm term}\) is genuine coefficient correlation:
\begin{enumerate}[label=(\roman*)]
 \item if \(\alpha=\beta\) and \(k\ne k'\), it contains
 \begin{equation}
  \mu(n_{\alpha,j,k})\mu(n_{\alpha,j,k'}),
  \qquad
  n_{\alpha,j,k}-n_{\alpha,j,k'}=(k-k')M;
  \label{eq:tpc39-terminal-two-Mobius}
 \end{equation}
 \item if \(\alpha\ne\beta\), the physical row provenance contributes
 the additional factors \(\mu(d_\alpha)\mu(d_\beta)\), producing the
 affine four--M\"obius quadratic form isolated in TPC--34 and TPC--38.
\end{enumerate}

\subsection{What minimal folding retains}

Let
\begin{equation}
 C_{k,k'}=\langle\mathscr Z_k,\mathscr Z_{k'}\rangle.
 \label{eq:tpc39-alias-covariance}
\end{equation}
The minimal twisted energy is exactly
\begin{align}
 E_{\rm fold}
 &=C_{0,0}
 \notag\\
 &\quad+
 \sum_{r=1}^{L}
 \bigl(
 C_{r,r}+C_{r-q,r-q}
 +2\operatorname{Re}C_{r,r-q}
 \bigr).
 \label{eq:tpc39-terminal-folded-covariance}
\end{align}
The diagonal entries in this formula still contain cross--row terms at a
fixed alias.  In addition, each folded pair has
\begin{equation}
 r-(r-q)=q,
 \qquad qM\in(B,B+M],
 \label{eq:tpc39-terminal-long-gap}
\end{equation}
so its cross term contains a long--gap M\"obius correlation.

The twist
\begin{equation}
 \ee{\nu(su-N_\alpha(j))/(qM)}
 =\omega^{\nu k}
 \label{eq:tpc39-terminal-physical-twist}
\end{equation}
is constant across all rows at fixed \(k\).  It therefore does not
oscillate the fixed--alias cross--row form.  The two members of a folded
pair have the same phase modulo \(q\), so that correlation also survives.
Most importantly, \(C_{0,0}=\|\mathscr R_0\|^2\) remains in
\eqref{eq:tpc39-terminal-folded-covariance} without alteration.

\subsection{Comparison with short progression variance}

The closest general theorem for scalar multiplicative functions in this
range is the variance estimate of
\citet[Theorem~1.5]{KlurmanMangerelTeravainen2023}.  After subtraction of
one most--pretentious character, it gives for almost every suitable
modulus \(m=o(x)\) a bound of schematic size
\begin{equation}
 \eps\,\varphi(m)\left(\frac xm\right)^2
 \label{eq:tpc39-KMT-scale}
\end{equation}
over the reduced residue classes; under GRH the exceptional--modulus
restriction is removed for typical moduli.  Specializing only the scale to
\(m=M\), \(x\asymp X\), and \(K=X/M\asymp K_B\), and using
\(\varphi(M)\asymp M\), the right side has order
\begin{equation}
 \eps MK^2,
 \label{eq:tpc39-KMT-endpoint-scale}
\end{equation}
whereas the scalar atomic diagonal has order \(MK\).  Thus this black box
would need \(\eps\ll K^{-1}X^{o(1)}\) to reach the diagonal scale.  At the
TPC endpoint, \(K=X^{1/400+o(1)}\), while the available savings are
logarithmic rather than this fixed power.  The M\"obius--specific
adaptation discussed by those authors likewise gives a logarithmic
saving, not \(K^{-1}\).

This comparison is deliberately limited.  The cited theorem averages
reduced residue classes, removes a rank--one pretentious component, and is
unconditional for almost all moduli.  It does not contain the physical
row mask, the row synthesis, or the four--M\"obius form in
\(\cO_{\rm term}\).  GRH changes the exceptional--modulus statement, not
the missing factor \(K\).  Hence no existing scalar progression theorem
is imported here as a proof of
\cref{ass:tpc39-minimal-twist-energy}.
